%% ========================================================================
%% Lampiran C: Topik Lanjutan
%% ========================================================================
%% STATUS: SCAFFOLD ONLY - lihat docs/authoring-guide.md
%% ========================================================================

\chapter{Topik Lanjutan}
\label{app:topik-lanjutan}

Simple POS, seperti yang dibangun sepanjang buku ini, sudah menjadi aplikasi
yang lengkap dan dapat dipakai. Namun ekosistem Laravel jauh lebih luas dari
yang tercakup pada dua belas bab inti. Lampiran ini merangkum secara singkat
tiga arah lanjutan yang lazim ditempuh pengembang Laravel setelah menguasai
materi buku ini, sebagai peta jalan, bukan sebagai materi yang diajarkan
penuh.

\section{Livewire}

% TODO: satu paragraf tentang Livewire sebagai alternatif membangun UI reaktif
% tanpa menulis JavaScript terpisah, dan bagaimana ia akan mengubah cara
% halaman kasir Simple POS dibangun dibanding pendekatan Blade + Alpine.js
% pada Bab 3.

\section{Inertia.js}

% TODO: satu paragraf tentang Inertia.js sebagai jembatan antara backend
% Laravel dan frontend berbasis komponen (Vue/React) tanpa membangun REST API
% terpisah, dan kapan pendekatan ini lebih cocok dibanding REST API murni
% pada Bab 9.

\section{Laravel Octane}

% TODO: satu paragraf tentang Laravel Octane untuk meningkatkan throughput
% aplikasi dengan menjaga aplikasi tetap termuat di memori antar-request
% (Swoole/RoadRunner), dan pertimbangan yang berubah dibanding deployment
% konvensional pada Bab 11.
